
\subsubsection{5.1 RQ1: RI Construct Validity (H-E1 — PASS)}

Both pre-registered gate conditions are met.

\textbf{SD(AdvGLUE_drop) = 0.1212 (threshold > 0.05; PASS).} The 95\% bootstrap CI is [0.093, 0.138], entirely above the threshold. The SD is 2.4 times the threshold value.

\textbf{R²_residualization = 0.529 (threshold < 0.80; PASS).} The 95\% bootstrap CI is [0.275, 0.721], entirely below the threshold. Capability and mean confidence together explain 52.9\% of adversarial drop variance, leaving 47.1\% as model-specific residual.

\begin{table}[htbp]
\centering
\begin{tabular}{lllll}
\toprule
Metric & Value & 95\% CI & Threshold & Status \\
\midrule
SD(AdvGLUE_drop) & 0.1212 & [0.093, 0.138] & > 0.05 & PASS \\
R²_residualization & 0.5285 & [0.275, 0.721] & < 0.80 & PASS \\
PC1 variance explained & 68.5\% & — & ≥ 70\% & WARNING \\
VIF (all covariates) & 1.000 & — & < 5.0 & PASS \\
\bottomrule
\end{tabular}
\end{table}

PC1 falls marginally below the 70\% variance threshold (68.5\%), introducing a minor residual capability confounding risk. VIF = 1.000 confirms that RI is orthogonal to PC1 by construction.

\subsubsection{5.2 RQ2: RI–ECE Partial Correlation (H-M1 — Significant but Inverted)}

\textbf{Primary result: ρ = −0.535, p = 0.0034, 95\% CI = [−0.782, −0.101], n = 30.}

The pre-registered prediction was ρ ≥ +0.4 (positive correlation: more adversarially fragile models predicted to have higher ECE). The observed correlation is −0.535 — statistically significant and in the opposite direction. Models with higher RI (greater adversarial fragility after capability control) tend to have lower ECE (better calibration on arc_challenge). The p-value of 0.0034 survives Holm-Bonferroni correction (α = 0.0125 for the primary prediction). The H-M1 gate therefore fails on two of three pre-registered conditions: the ρ sign condition (expected positive, observed negative) and the family sign consistency condition (expected ≥2/3 positive, observed 1/3 positive).

The gate failure on H-M1, which was designated MUST_WORK, blocked execution of H-M2 (HaluEval), H-M3 (HarmBench), and H-M4 (OVI-GSM8K). Those sub-hypotheses remain not executed.

\textbf{Robustness checks:}

Three outliers were identified by Cook's distance: meta-llama/Meta-Llama-3-70B, google/gemma-7b-it, and stabilityai/stablelm-zephyr-3b. Excluding these three, ρ = −0.498, p = 0.008 — the direction and significance are preserved.

Fisher z-test for interaction with capability level: z = −0.561, p = 0.575 (not significant). The anticorrelation does not vary significantly with capability level.

Baseline comparisons:
\begin{itemize}
\item ρ(PC1, ECE) = −0.511, p = 0.0039 (capability-only predictor)
\item ρ(raw AdvGLUE_drop, ECE) = −0.418, p = 0.0213 (uncontrolled)
\end{itemize}

The similarity between ρ(PC1, ECE) = −0.511 and ρ(RI, ECE | PC1) = −0.535 indicates that RI adds limited unique predictive signal for ECE beyond what capability (PC1) already captures.

\begin{table}[htbp]
\centering
\begin{tabular}{llll}
\toprule
Condition & Threshold & Actual & Status \\
\midrule
Spearman partial ρ(RI, ECE) & ≥ +0.4 & −0.535 & FAIL (inverted sign) \\
Holm-corrected p-value & < 0.05 & 0.0034 & PASS \\
Consistent positive families & ≥ 2/3 & 1/3 (Qwen only) & FAIL \\
\bottomrule
\end{tabular}
\end{table}

\subsubsection{5.3 Per-Family Analysis}

\begin{table}[htbp]
\centering
\begin{tabular}{llllll}
\toprule
Family & n & ρ(RI, ECE) & p (raw) & p (Holm) & Direction \\
\midrule
LLaMA & 9 & −0.244 & 0.599 & 1.000 & Negative \\
Mistral & 6 & −0.827 & 0.173 & 0.519 & Negative \\
Qwen & 6 & +0.364 & 0.636 & 1.000 & Positive \\
\bottomrule
\end{tabular}
\end{table}

Two of three families show negative RI–ECE correlations. No within-family result is statistically significant after Holm correction. Mistral shows the strongest point estimate (ρ = −0.827), but with n = 6 this is underpowered (raw p = 0.173, Holm p = 0.519). Qwen shows a positive sign that differs from the aggregate direction; with n = 6 and p = 0.636, this cannot be interpreted reliably. Per-family results are treated as exploratory.

\subsubsection{5.4 Reliability Diagram Analysis}

Average reliability diagrams stratified by RI quartile show that models in the highest RI quartile (Q4, most adversarially fragile) have calibration curves closest to the perfect diagonal on arc_challenge, while models in the lowest RI quartile (Q1, most robust) show curves further from the diagonal, indicating greater overconfidence. This qualitative pattern is consistent with the statistical result.

